\chapter[Band 32: Kommutative idempotente Halbgruppen]{Band 32 auf einen Blick}
\noindent\textbf{Kommutative idempotente Halbgruppen}\par\medskip
Die beiden vorherigen Strukturaxiomatiken werden gemeinsam gefordert.
Ihre ordnungstheoretische Lesart wird in Band 45 entwickelt.

\begin{BandUebersicht}
\textbf{Gemeinsame Struktur}
\UebFormel{\begin{aligned}
&\CommutativeIdempotentSemigroupAlg A\star\
 \leftrightarrow\\
&\quad\CommutativeSemigroupAlg A\star\land\\
&\quad\IdempotentSemigroupAlg A\star.
\end{aligned}}
\UebQuellen{\FormulaRefAuto{CommutativeIdempotentSemigroupDef}[def]}
&
\textbf{Vereinigung erfüllt beide Axiome}
\UebFormel{\begin{aligned}
&\CommutativeIdempotentSemigroupAlg{\powerset(M)}{\cup},\\
&X,Y\in\powerset(M)\ \vdash\\
&\quad X\cup Y=Y\cup X,\qquad X\cup X=X.
\end{aligned}}
\UebQuellen{\FormulaRefAuto{PowerSetUnionCommutativeIdempotentSemigroup}[thm];
\FormulaRefAuto{CommutativeIdempotentSemigroupCommutativityAxiom}[ax];
\FormulaRefAuto{CommutativeIdempotentSemigroupIdempotenceAxiom}[ax]} \\
\addlinespace[.8ex]

\textbf{Homomorphismen und Isomorphismen}
Für zwei Strukturen derselben Klasse bedeutet ein Homomorphismus
\(H(f)\):\UebFormel{\begin{aligned}
&f:A\to B,\\
&x,y\in A\ \vdash\\
&\quad f(x\star y)=f(x)\diamond f(y).
\end{aligned}}Ein Isomorphismus \(I(f)\) ist ein bijektiver Homomorphismus:\UebFormel{\begin{aligned}
&I(f)\ \leftrightarrow H(f)\land f:A\bij B.
\end{aligned}}
\UebQuellen{\FormulaRefAuto{CommutativeIdempotentSemigroupHomDef}[def];
\FormulaRefAuto{CommutativeIdempotentSemigroupIsoDef}[def]}
&
\textbf{Komposition erhält die Struktur}
Für Homomorphismen \(f:A\to B\), \(g:B\to C\) zwischen
Strukturen dieser Klasse ist \(g\circ f\) wieder ein Homomorphismus.
Mit \(\bullet\) als Operation auf \(C\) gilt für \(x,y\in A\):\UebFormel{\begin{aligned}
&(g\circ f)(x\star y)\\
&\quad=(g\circ f)(x)\mathbin{\bullet}(g\circ f)(y).
\end{aligned}}Sind \(f,g\) Isomorphismen, so ist auch \(g\circ f\) ein Isomorphismus.
\UebQuellen{\FormulaRefAuto{SemigroupHomComposition}[thm];
\FormulaRefAuto{SemigroupIsoComposition}[thm]} \\
\addlinespace[.8ex]
\end{BandUebersicht}

